\chapter{Building Blocks of the Transport}
\label{ch:transport}

\epigraph{``Show me the money!''}{Rod Tidwell, \emph{Jerry Maguire} (1996)}

Rod's point was not greed. It was that he had spent years being told about
value in the abstract and wanted to see the actual figure. This chapter is that
demand, applied to MCP. We are going to look at literal bytes on a literal
connection, because every claim about performance from here on rests on knowing
what is genuinely on the wire and what you imagined was there.

There are two transports. There used to be three.

\section{JSON-RPC, and the two things MCP adds}

The framing is JSON-RPC 2.0, unmodified, which was a good decision made by
people resisting the urge to design something. It has three message types.

\textbf{Requests} have an id and expect a response:

\begin{wire}
{ "jsonrpc": "2.0", "id": 1, "method": "tools/call", "params": { ... } }
\end{wire}

\textbf{Results} come back with the same id:

\begin{wire}
{ "jsonrpc": "2.0", "id": 1, "result": { "resultType": "complete", ... } }
\end{wire}

\textbf{Notifications} have no id and get no reply:

\begin{wire}
{ "jsonrpc": "2.0", "method": "notifications/progress", "params": { ... } }
\end{wire}

MCP adds exactly two rules on top:

\begin{enumerate}[leftmargin=1.6em, itemsep=3pt]
\item \texttt{result.resultType} is required. Chapter~\ref{ch:architecture}
covered why: it is the hinge that lets MRTR and Tasks exist without changing
anything else.
\item Request ids must not be \texttt{null}, and must be unique among requests
you have issued and not yet had answered. Base JSON-RPC allows a null id. MCP
does not, because null ids make correlation ambiguous on a shared channel.
\end{enumerate}

That is the whole extension. Everything interesting is in the payloads.

\subsection{Error codes, and the range you must not touch}

JSON-RPC reserves \texttt{-32768} to \texttt{-32000}. Within that, it leaves
\texttt{-32000} to \texttt{-32099} for implementations, and \texttt{2026-07-28}
finally partitioned that space, which needed doing.

\begin{table}[htbp]
\centering
\small
\begin{tabularx}{\textwidth}{@{}llX@{}}
\toprule
\thead{Range} & \thead{Status} & \thead{Rule} \\
\midrule
\texttt{-32700}, \texttt{-32600} to \texttt{-32603}
  & JSON-RPC standard & Parse, request, method, params, internal \\
\texttt{-32000} to \texttt{-32019}
  & Legacy & Allocated before the policy existed. Do not allocate here, and do
    not assume any meaning \\
\texttt{-32020} to \texttt{-32099}
  & \textbf{Reserved for MCP} & Only the specification allocates here. Never
    emit an undefined code from this range \\
Outside \texttt{-32768..-32000}
  & Yours & Put your application errors here \\
\bottomrule
\end{tabularx}
\caption{The error-code partition. The middle row is the one people violate,
usually by picking a ``nice round'' number like \texttt{-32050}.}
\label{tab:error-ranges}
\end{table}

Three codes are currently defined in the reserved range, and each tells the
client something actionable:

\begin{table}[htbp]
\centering
\small
\begin{tabularx}{\textwidth}{@{}llX@{}}
\toprule
\thead{Code} & \thead{Name} & \thead{What the client should do} \\
\midrule
\texttt{-32020} & \texttt{HeaderMismatch}
  & Refetch \texttt{tools/list} (the schema may have changed), then retry with
    correct headers \\
\texttt{-32021} & \texttt{MissingRequiredClientCapability}
  & Read \texttt{data.requiredCapabilities}; either declare them or stop \\
\texttt{-32022} & \texttt{UnsupportedProtocolVersion}
  & Pick a version from \texttt{data.supported} and retry \\
\bottomrule
\end{tabularx}
\caption{Every one of these carries enough information for the client to
recover automatically. That is not an accident, it is the design goal.}
\label{tab:mcp-errors}
\end{table}

Two codes are retired and must never be re-emitted: \texttt{-32002}, which meant
``resource not found'' until this revision replaced it with \texttt{-32602}, and
\texttt{-32042}, which meant ``URL elicitation required'' for exactly one
revision. Clients should still \emph{accept} \texttt{-32002} from older servers.
Servers on this revision must not send it. Meridian has a test for precisely
this, because it is the kind of thing that rots quietly:

\begin{py}
def test_retired_codes_are_never_emitted(self):
    server = tiny_server()
    seen = set()
    for method, params in [("resources/read", {"uri": "x://y"}),
                           ("tools/call", {"name": "nope"}),
                           ("prompts/get", {"name": "nope"})]:
        resp = server.handle(request(method, params))
        seen.add(resp["error"]["code"])
    self.assertNotIn(-32002, seen)
    self.assertNotIn(-32042, seen)
\end{py}

\section{Streamable HTTP}

One endpoint. POST only. Every JSON-RPC message is its own HTTP request.

That sentence is doing a lot of work, because it describes something quite
different from what the same transport looked like a year ago. Here is what
2026-07-28 deleted:

\begin{legacybox}{What Streamable HTTP used to have}
Protocol versions \texttt{2025-03-26} through \texttt{2025-11-25} used the same
transport name and a materially different shape:

\begin{itemize}[leftmargin=1.4em, itemsep=2pt]
\item \texttt{Mcp-Session-Id}, assigned by the server, terminated with HTTP DELETE
\item A standalone \texttt{GET} that opened an SSE stream for server-initiated messages
\item Servers sending JSON-RPC \emph{requests} on those streams
\item Stream resumability via \texttt{Last-Event-ID} and SSE event ids
\end{itemize}

None of it survives. A server on this revision should answer \texttt{405} to
GET and DELETE, ignore \texttt{Mcp-Session-Id} without echoing it, and ignore
\texttt{Last-Event-ID}.
\end{legacybox}

\subsection{The three shapes of traffic}

\bookfig[0.68]{http-flows}{Everything that happens on a Streamable HTTP
connection. Three shapes, and the third one is still request/response no matter
how long it stays open.}{http-flows}

The elegance is in the third one. A long-lived notification stream is not a
special connection type with its own lifecycle. It is an ordinary POST whose
response happens to be an SSE stream that stays open. Its state belongs to the
request, not to the connection, which is why it survives the deletion of
sessions without needing a replacement concept.

\subsection{Required headers, and why they exist}
\label{sec:required-headers}

Here is a genuinely new obligation. Every POST must carry these:

\begin{http}
POST /mcp HTTP/1.1
Content-Type: application/json
Accept: application/json, text/event-stream
MCP-Protocol-Version: 2026-07-28
Mcp-Method: tools/call
Mcp-Name: assess_account_risk

{
  "jsonrpc": "2.0",
  "id": 1,
  "method": "tools/call",
  "params": {
    "name": "assess_account_risk",
    "arguments": { "accountId": "ACC-1042" },
    "_meta": {
      "io.modelcontextprotocol/protocolVersion": "2026-07-28",
      "io.modelcontextprotocol/clientCapabilities": {}
    }
  }
}
\end{http}

\texttt{Mcp-Method} mirrors \texttt{method}. \texttt{Mcp-Name} mirrors
\texttt{params.name} (or \texttt{params.uri} for \texttt{resources/read}), and
is required for \texttt{tools/call}, \texttt{prompts/get}, and
\texttt{resources/read}.

The point is that an intermediary can now do its job without parsing JSON:

\begin{itemize}[leftmargin=1.6em, itemsep=3pt]
\item A load balancer routes \texttt{tools/call} to a compute pool and
\texttt{tools/list} to a cache tier.
\item A WAF applies a strict policy to one tool name and a relaxed one to
another.
\item Rate limiting works per method rather than per connection, which is the
only granularity that means anything now that connections are meaningless.
\item Your access logs become useful without a JSON parser in the log pipeline.
\end{itemize}

\begin{dangerbox}{Two sources of truth is a smuggling primitive}
Mirroring creates an obvious hazard. If a gateway routes on \texttt{Mcp-Name:
read\_only\_tool} and the body says \texttt{"name": "wire\_transfer"}, the
gateway has approved one thing and the server has executed another.

So validation is mandatory, not advisory. Any server that processes the body
\textbf{must} reject a mismatch with \texttt{400} and \texttt{-32020}.
\end{dangerbox}

\begin{py}
def _validate_headers(self, headers, message: dict) -> None:
    """Reject any disagreement between headers and body."""
    method = message.get("method", "")
    params = message.get("params") or {}

    version_header = headers.get(HDR_VERSION)
    if version_header is None:
        raise errors.HeaderMismatch(f"Missing required header {HDR_VERSION}")
    body_version = (params.get("_meta") or {}).get(KEY_PROTOCOL_VERSION)
    if body_version is not None and version_header != body_version:
        raise errors.HeaderMismatch(
            f"{HDR_VERSION} header {version_header!r} does not match "
            f"body value {body_version!r}")

    method_header = headers.get(HDR_METHOD)
    if method_header != method:
        raise errors.HeaderMismatch(
            f"{HDR_METHOD} header {method_header!r} does not match "
            f"body method {method!r}")

    source = NAME_SOURCE.get(method)
    if source:
        name_header = headers.get(HDR_NAME)
        if decode_header_value(name_header) != params.get(source):
            raise errors.HeaderMismatch(
                f"{HDR_NAME} header does not match body {source}")
\end{py}

Meridian tests each failure mode separately, because ``we validate headers'' is
the kind of claim that is true until somebody refactors:

\begin{py}
def test_method_header_disagreeing_with_body_is_rejected(self):
    """The request-smuggling case. A gateway routing on the header must
    never be able to disagree with the server executing on the body."""
    status, body = self.raw_post(
        {"jsonrpc": "2.0", "id": 1, "method": "tools/call",
         "params": {"name": "echo", "arguments": {"text": "hi"},
                    "_meta": meta()}},
        self.good_headers("tools/list", "echo"))
    self.assertEqual(status, 400)
    self.assertEqual(body["error"]["code"], errors.HEADER_MISMATCH)
\end{py}

\subsection{Mirroring tool parameters into headers}

Servers can go further and expose specific tool arguments as headers, using an
\texttt{x-mcp-header} annotation in the input schema:

\begin{wire}
{
  "name": "assess_account_risk",
  "inputSchema": {
    "type": "object",
    "properties": {
      "accountId": { "type": "string" },
      "region": {
        "type": "string",
        "enum": ["us-east", "us-west", "eu-central"],
        "x-mcp-header": "Region"
      }
    },
    "required": ["accountId"]
  }
}
\end{wire}

A call with \texttt{"region": "eu-central"} now carries
\texttt{Mcp-Param-Region: eu-central}, and your gateway can route EU traffic to
EU infrastructure without reading a byte of the body. For a data-residency
requirement that is genuinely valuable.

The constraints are strict, and each one closes a specific hole:

\begin{table}[htbp]
\centering
\small
\begin{tabularx}{\textwidth}{@{}lX@{}}
\toprule
\thead{Constraint} & \thead{Why} \\
\midrule
Primitive types only (string, integer, boolean; \emph{not} number)
  & Floats do not have a canonical string form, so header and body could
    disagree while both being ``correct'' \\
Case-insensitively unique within the schema
  & Header names are case-insensitive, so two annotations differing only in
    case would collide \\
No control characters, including CR and LF
  & Header injection \\
Statically reachable from the root, through \texttt{properties} keys only
  & An intermediary must know from the schema alone which headers a call will
    carry. No arrays, no \texttt{oneOf}, no \texttt{\$ref} \\
\bottomrule
\end{tabularx}
\caption{The \texttt{x-mcp-header} rules. The last one is the subtle one, and
Meridian validates it recursively.}
\label{tab:x-mcp-header}
\end{table}

A client on Streamable HTTP must \emph{reject} a tool whose annotations violate
these rules, and rejection means excluding that one tool from
\texttt{tools/list} rather than failing the whole catalogue. One bad tool
should not take down the other forty-nine.

\subsection{Encoding values safely}

RFC 9110 allows visible ASCII, space, and tab in header values. Anything else
gets wrapped in a sentinel:

\begin{table}[htbp]
\centering
\small
\begin{tabular}{@{}lll@{}}
\toprule
\thead{Value} & \thead{Why encoded} & \thead{Header} \\
\midrule
\texttt{us-west1}       & plain ASCII        & \texttt{us-west1} \\
\texttt{caf\'e}          & non-ASCII          & \texttt{=?base64?Y2Fmw6k=?=} \\
\texttt{"\ padded\ "}   & surrounding spaces & \texttt{=?base64?IHBhZGRlZCA=?=} \\
\texttt{line1\textbackslash nline2} & newline & \texttt{=?base64?bGluZTEKbGluZTI=?=} \\
\texttt{=?base64?literal?=} & looks like the sentinel & \texttt{=?base64?PT9iYXNlNjQ/bGl0ZXJhbD89?=} \\
\bottomrule
\end{tabular}
\caption{Header value encoding. The last row is the one that separates a
correct implementation from a plausible one.}
\label{tab:header-encoding}
\end{table}

That last row deserves a moment. A plain ASCII value that happens to look like
the sentinel must \emph{also} be encoded, otherwise a decoder turns
\texttt{"=?base64?literal?="} into something the sender never wrote. It is a
tiny rule and exactly the sort of thing that produces a security advisory two
years later.

\begin{py}
def encode_header_value(value) -> str:
    if isinstance(value, bool):
        text = "true" if value else "false"
    else:
        text = str(value)

    needs_encoding = (
        not text.isascii()
        or any(ord(c) < 0x20 or ord(c) == 0x7F for c in text)
        or text != text.strip()
        or (text.startswith(B64_PREFIX) and text.endswith(B64_SUFFIX))
    )
    if not needs_encoding:
        return text
    blob = base64.b64encode(text.encode("utf-8")).decode("ascii")
    return f"{B64_PREFIX}{blob}{B64_SUFFIX}"
\end{py}

\subsection{No resumability, and what that means for you}

A broken response stream loses the request. There is no \texttt{Last-Event-ID},
no redelivery, no ``resume from event 14''. The client re-issues with a
\emph{new} request id and starts over.

This is a real simplification, and it moves a burden onto you.

\keyidea{Design your servers to be idempotent, because retries are now a normal
part of operation rather than an exception. If re-issuing \texttt{transfer\_funds}
twice moves money twice, the transport will eventually find out for you.}

The usual mechanism is an idempotency key supplied by the caller and honoured
for a window on the server. It is not a protocol feature; it is an ordinary tool
argument, and Chapter~\ref{ch:building-servers} builds one.

\subsection{Cancellation}

On HTTP, closing the response stream \emph{is} cancellation. That is the whole
mechanism. Because each request has its own stream, a transport-level disconnect
is unambiguous, so no \texttt{notifications/cancelled} is sent or expected.

The server must stop as soon as practical and must send nothing further for that
request. In Meridian this falls out of the write path:

\begin{py}
def write(msg: dict) -> None:
    with lock:
        try:
            h.wfile.write(jsonrpc.sse_event(msg))
            h.wfile.flush()
        except (BrokenPipeError, ConnectionResetError):
            raise _ClientGone()
\end{py}

A client hanging up mid-stream is normal here, not an incident. Which means the
stock Python HTTP server's habit of printing a traceback for every reset
connection turns routine operation into alarming log noise, so Meridian
suppresses it explicitly. Small thing. Saves somebody a 3 a.m. investigation.

\subsection{One header that is pure profit}

\begin{http}
X-Accel-Buffering: no
\end{http}

Reverse proxies buffer responses by default. Without this header nginx will
happily accumulate your progress notifications and deliver them in one lump
after the final response, at which point you have all the complexity of
streaming and none of the benefit. Send it on every SSE response.

\section{stdio}

The other transport is a subprocess and two pipes. It is much simpler and it
has sharper edges.

\begin{itemize}[leftmargin=1.6em, itemsep=3pt]
\item The client launches the server as a child process.
\item Messages go over \texttt{stdin} and \texttt{stdout}, one per line.
\item Messages must not contain embedded newlines.
\item \texttt{stderr} is for logging, and the client may capture, forward, or
ignore it.
\item The server must write \textbf{nothing} to \texttt{stdout} that is not a
valid MCP message.
\end{itemize}

That last rule is the one that will bite you, and it will not be your code that
breaks it. It will be a library that prints a deprecation warning to stdout on
import. One stray \texttt{print()} corrupts the stream, and the client sees a
parse error it cannot attribute to anything.

Meridian encodes with no whitespace and asserts the invariant in a test:

\begin{py}
def test_encoded_messages_contain_no_newlines(self):
    """stdio framing depends on this and nothing enforces it but us."""
    wire = encode({"jsonrpc": "2.0", "id": 1,
                   "result": {"text": "line one\nline two"}})
    self.assertNotIn("\n", wire)
\end{py}

\subsection{A connection is not a conversation}

The specification is unusually direct about this, and it is worth quoting the
consequence: clients may interleave unrelated requests on the same pipe, and a
server must not treat process identity as a proxy for conversation continuity.

Practically: do not put a user id in a module-level variable at startup. Do not
cache ``the current account'' between calls. The next message on that pipe may
belong to an entirely different task.

\subsection{Cold start is the number that decides your topology}

Because every message shares one channel and the process must be spawned, stdio
has a cost HTTP does not: process startup.

\begin{measurebox}{stdio cold start versus warm call}
\begin{tabular}{@{}lr@{}}
\toprule
\thead{} & \thead{Time} \\
\midrule
Spawn to first usable response  & 44.6\,ms \\
Warm call, same server          & 0.060\,ms \\
\midrule
Calls needed to amortise a cold start & \textasciitilde{}750 \\
\bottomrule
\end{tabular}

\vspace{6pt}
\footnotesize
Python 3.14 on arm64. Interpreter boot plus imports. A Node server is
comparable; a static binary is far better.
\end{measurebox}

Forty-five milliseconds sounds small until you notice it is 750 times a warm
call. Spawn-per-call is therefore indefensible for anything but a one-shot CLI.
Keep a warm pool, and Chapter~\ref{ch:optimizing} shows how.

\subsection{Shutdown, done properly}

Close stdin. Wait. Escalate only if you must.

\begin{py}
def close(self, timeout: float = 5.0) -> None:
    try:
        self._proc.stdin.close()
    except Exception:
        pass
    try:
        self._proc.wait(timeout=timeout)
    except subprocess.TimeoutExpired:
        # Escalate only if the server ignored the EOF on stdin.
        self._proc.terminate()
        try:
            self._proc.wait(timeout=2.0)
        except subprocess.TimeoutExpired:
            self._proc.kill()
\end{py}

Closing stdin is the primary graceful-shutdown signal and the only portable one,
so servers should exit promptly on EOF. If yours does not, every client that
talks to it will eventually SIGKILL it, and your cleanup handlers will never run.
Meridian tests that its servers exit cleanly rather than needing to be killed.

\subsection{Cancellation on stdio}

Different from HTTP, and for a good reason: there is no per-request stream to
close, so cancellation is an explicit notification.

\begin{wire}
{ "jsonrpc": "2.0", "method": "notifications/cancelled",
  "params": { "requestId": 42 } }
\end{wire}

The server should stop work and must send nothing further for that request,
including its response.

\begin{notebox}{The asymmetry is not sloppiness}
It is the transport binding doing its job. On HTTP the disconnect is
unambiguous, so a message would be redundant. On stdio the channel is shared, so
a message is the only signal available. Same semantics, different mechanism,
which is exactly what a transport binding is for.
\end{notebox}

\section{Comparing the two, with numbers}

\bookfig[0.98]{transport-cost}{Left: protocol overhead with server execution
subtracted out. Right: the same numbers next to actual network distance, on a
log scale, which is the honest context.}{transport-cost}

\begin{measurebox}{Same call, three transports}
\begin{tabular}{@{}lrrr@{}}
\toprule
\thead{Transport} & \thead{Mean} & \thead{p95} & \thead{Overhead} \\
\midrule
In-process (control) & 0.031\,ms & 0.034\,ms & baseline \\
stdio                & 0.060\,ms & 0.065\,ms & +0.029\,ms \\
Streamable HTTP      & 0.161\,ms & 0.203\,ms & +0.130\,ms \\
\bottomrule
\end{tabular}

\vspace{6pt}
\footnotesize
The control still serialises and deserialises, so this compares transport
mechanics rather than flattering the in-process case.
\end{measurebox}

Read that right-hand chart carefully, because it is the point of the section.
Streamable HTTP costs 0.13\,ms more than in-process on loopback. A same-region
network hop costs around 1.2\,ms. A cross-region hop costs around 68\,ms.

\keyidea{Transport overhead is not your problem. Distance is your problem, and
round trips are how you pay for distance. Choose your transport on operational
grounds, then spend your optimisation effort on the number of trips.}

\begin{table}[htbp]
\centering
\small
\begin{tabularx}{\textwidth}{@{}lXX@{}}
\toprule
& \thead{stdio} & \thead{Streamable HTTP} \\
\midrule
Deployment      & Subprocess, same machine & Network service \\
Isolation       & Process boundary & Process, container, or host \\
Startup         & 44.6\,ms per spawn & Connection setup, then reused \\
Concurrency     & One shared channel & One stream per request \\
Cancellation    & \texttt{notifications/cancelled} & Close the stream \\
Auth            & Environment (credentials from the parent) & OAuth 2.1 (\chapref{authorization}) \\
Scaling         & One process per host & Round-robin, any replica \\
Best for        & Local tools, filesystems, developer machines & Everything shared or remote \\
\bottomrule
\end{tabularx}
\caption{Choosing. Appendix~\ref{ch:transport-matrix} has the long version.}
\label{tab:transport-choice}
\end{table}

\section{Topologies}

The stateless core changes what is possible here, so it is worth drawing.

\bookfig[0.98]{topologies}{Three shapes. The middle one is the one that got
dramatically easier in 2026-07-28, because round-robin now works without a
session store behind it.}{topologies}

\textbf{Sidecar} is stdio, one server process per host, no network. Simple
isolation, and you pay the cold start.

\textbf{Gateway} is the interesting one. Before this revision, N hosts behind
one gateway meant sticky sessions, because request two had to reach whichever
replica handled request one. Now it does not, so the gateway is a plain
round-robin load balancer with a cache in front of it, and the routing headers
from \secref{required-headers} let it make decisions without parsing bodies.

\textbf{Mesh} is agents talking to agents, where a node is a host on one side
and a server on the other. Chapter~\ref{ch:multi-agent} builds this and spends
most of its time on the failure modes, because a mesh with no depth limit is a
fork bomb with a nicer API.

\section{Meridian's transport decision}

Meridian runs four servers. Here is what each got and why, which is a more
useful format than a general recommendation.

\begin{table}[htbp]
\centering
\small
\begin{tabularx}{\textwidth}{@{}llX@{}}
\toprule
\thead{Server} & \thead{Transport} & \thead{Reason} \\
\midrule
\texttt{risk} & HTTP
  & Shared across three business units, needs OAuth, must scale horizontally \\
\texttt{compliance} & HTTP
  & Audit trail must be centralised, and it sits behind an API gateway \\
\texttt{fraud} & HTTP
  & Called in parallel with risk, so it wants connection reuse \\
\texttt{marketdata} & HTTP
  & \texttt{cacheScope: public} results, so a shared gateway cache is the entire
    point \\
\midrule
Local dev & stdio & No auth, no ports, and \texttt{.mcp.json} just works \\
\bottomrule
\end{tabularx}
\caption{Meridian's choices. Everything shared is HTTP; stdio is a development
convenience and a good one.}
\label{tab:meridian-transports}
\end{table}

The \texttt{marketdata} row is worth noticing. It is the only server whose
results are identical for every caller, which is what makes
\texttt{cacheScope: public} honest, which is what makes a shared gateway cache
legal. Chapter~\ref{ch:resources} measures that: 577 of 600 reads avoided, at a
96.2\,\% hit rate.

\section{What to remember}

\begin{itemize}[leftmargin=1.6em, itemsep=3pt]
\item JSON-RPC 2.0, plus required \texttt{resultType} and non-null ids.
\item Never allocate error codes in \texttt{-32020} to \texttt{-32099}. Never
re-emit \texttt{-32002} or \texttt{-32042}.
\item Streamable HTTP is POST-only. GET, DELETE, sessions, and resumability are
all gone.
\item \texttt{MCP-Protocol-Version}, \texttt{Mcp-Method}, and \texttt{Mcp-Name}
are mandatory, and a mismatch between header and body must be rejected.
\item No resumability means retries are routine. Make your servers idempotent.
\item Cancellation is a closed stream on HTTP and a notification on stdio.
\item stdio cold start is 44.6\,ms, around 750 warm calls. Use a warm pool.
\item Transport overhead is fractions of a millisecond. Distance and round trips
are what cost you.
\end{itemize}

Next: authorization, which is a security chapter with a performance chapter
hiding inside it.
